% documentation.tex
% Compilar con: pdflatex -output-directory=out documentation.tex

\documentclass[
  % -------------------------
  % Curso
  % -------------------------
  coursetitle={Manual de uso de \texttt{lbustio\_lecture\_notes}},
  coursehours={96},
  coursedate={Verano 2026},
  doctype={Documentaci\'on de la clase},
  otherlangs={english,french},
  % -------------------------
  % Autor
  % -------------------------
  authorname={L\'azaro},
  authorsurname={Bustio-Mart\'inez},
  authortitle={PhD},
  role={Profesor Investigador},
  email={lbustio@inaoe.mx},
  phone={5559504000},
  ext={7459},
  % -------------------------
  % Instituci\'on
  % -------------------------
  institution={Instituto Nacional de Astrof\'isica, \'Optica y Electr\'onica (INAOE)},
  department={Maestr\'ia en Ciencias en Tecnolog\'ias de Seguridad},
  unit={Coordinaci\'on de Ciberseguridad},
  address={Luis Enrique Erro 1, Sta. Ma. Tonantzintla, San Andr\'es Cholula, CP 72840, Puebla, M\'exico},
  % -------------------------
  % M\'argenes
  % -------------------------
  margin={0.5in}
]{lbustio_lecture_notes}

\usepackage{lipsum}

\begin{document}

\IberoCourseInfo

\section{Prop\'osito}
Este documento describe la estructura actual de la clase \texttt{lbustio\_lecture\_notes}. La clase separa los datos del documento de su presentaci\'on visual: primero se declaran datos de autor, instituci\'on y curso; despu\'es la clase usa esos datos para construir el encabezado, la informaci\'on del curso y el pie de p\'agina.

\section{Modelo de datos}
La clase define tres grupos de datos. Estos grupos no son secciones visuales; son solamente la fuente de informaci\'on.

Si un campo queda vac\'io, la clase no lo imprime en el header, el footer ni Course Info. Esto tambi\'en evita separadores sueltos: no aparecen comas, guiones, puntos ni l\'ineas vac\'ias para datos que no fueron declarados.

\subsection{Autor}
El autor representa a la persona responsable del documento. Se declara con:

\begin{verbatim}
authorname={L\'azaro},
authorsurname={Bustio-Mart\'inez},
authortitle={PhD},
role={Profesor Investigador},
email={lbustio@inaoe.mx},
phone={5559504000},
ext={7459},
\end{verbatim}

\subsection{Instituci\'on}
La instituci\'on representa la adscripci\'on del autor. Se declara con:

\begin{verbatim}
institution={Instituto Nacional de Astrof\'isica, \'Optica y Electr\'onica (INAOE)},
department={Maestr\'ia en Ciencias en Tecnolog\'ias de Seguridad},
unit={Coordinaci\'on de Ciberseguridad},
address={Luis Enrique Erro 1, Sta. Ma. Tonantzintla, San Andr\'es Cholula,
CP 72840, Puebla, M\'exico},
\end{verbatim}

\subsection{Curso}
El curso contiene los datos propios del material acad\'emico:

\begin{verbatim}
coursetitle={C\'omputo Nube y Tecnolog\'ias Emergentes},
coursehours={96},
coursedate={Verano 2026},
doctype={Gu\'ia del curso},
\end{verbatim}

\section{Secciones visuales}
A partir de los tres grupos de datos, la clase genera tres secciones visuales.

\subsection{Header}
El encabezado aparece en todas las p\'aginas. Usa el mismo estilo que el footer: texto peque\~no, gris oscuro, alineado a la derecha y separado del cuerpo por una l\'inea gris. Contiene:

\begin{itemize}
  \item Nombre y apellidos del autor, con su t\'itulo.
  \item Rol del autor.
  \item Email, tel\'efono y extensi\'on.
\end{itemize}

\subsection{Course Info}
La secci\'on de informaci\'on del curso se inserta con:

\begin{verbatim}
\IberoCourseInfo
\end{verbatim}

Por compatibilidad, \verb|\IberoCover| apunta al mismo macro. Esta secci\'on contiene:

\begin{itemize}
  \item T\'itulo del curso entre dos l\'ineas azules.
  \item Instituci\'on, departamento y unidad.
  \item Autor y t\'itulo del autor.
  \item Email del autor.
  \item Fecha del curso.
  \item Cantidad de horas.
  \item Tipo de documento.
\end{itemize}

\subsection{Footer}
El pie de p\'agina aparece en todas las p\'aginas. Usa el mismo tama\~no, color y l\'inea separadora que el header. Contiene:

\begin{itemize}
  \item Instituci\'on, departamento y unidad, alineados a la derecha.
  \item Direcci\'on de la instituci\'on, alineada a la derecha.
  \item N\'umero de p\'agina centrado.
\end{itemize}

\section{Ejemplo m\'inimo}
El siguiente ejemplo muestra la estructura recomendada:

\begin{verbatim}
\documentclass[
  coursetitle={C\'omputo Nube y Tecnolog\'ias Emergentes},
  coursehours={96},
  coursedate={Verano 2026},
  doctype={Gu\'ia del curso},

  authorname={L\'azaro},
  authorsurname={Bustio-Mart\'inez},
  authortitle={PhD},
  role={Profesor Investigador},
  email={lbustio@inaoe.mx},
  phone={5559504000},
  ext={7459},

  institution={Instituto Nacional de Astrof\'isica, \'Optica y Electr\'onica (INAOE)},
  department={Maestr\'ia en Ciencias en Tecnolog\'ias de Seguridad},
  unit={Coordinaci\'on de Ciberseguridad},
  address={Luis Enrique Erro 1, Sta. Ma. Tonantzintla, San Andr\'es Cholula,
  CP 72840, Puebla, M\'exico},

  margin={0.5in}
]{lbustio_lecture_notes}

\begin{document}
\IberoCourseInfo

\section{Contenido}
...
\end{document}
\end{verbatim}

\section{Notas de compatibilidad}
La clase ya no usa una entidad de datos separada llamada \texttt{Instructor}. Tampoco hay herencia entre instructor y autor. La informaci\'on de Course Info sale directamente de los datos del autor, la instituci\'on y el curso.

Si un documento viejo usa llaves como \texttt{instructor}, \texttt{instructortitle} o \texttt{instructoremail}, debe migrarse al modelo nuevo. Tambi\'en se recomienda usar \texttt{institution} en lugar de \texttt{university}; \texttt{university} queda solamente como respaldo interno cuando \texttt{institution} no se declara.

\section{Diagn\'ostico r\'apido}
Si aparece \texttt{Runaway argument?} cerca de \verb|\documentclass|, casi siempre falta cerrar una llave o un corchete en las opciones de la clase.

\begin{itemize}
  \item Verifica que cada \verb|key={...}| cierre con \verb|}|.
  \item Verifica que la lista completa cierre con \verb|]|.
  \item Evita dejar una coma final justo antes de \verb|]{lbustio_lecture_notes}|.
\end{itemize}

\section{Contenido de prueba}
Este contenido solo sirve para verificar que el header y el footer se repiten correctamente en varias p\'aginas.

\subsection{Secci\'on 1}
\lipsum[1-4]

\subsection{Secci\'on 2}
\lipsum[5-8]

\subsection{Secci\'on 3}
\lipsum[9-12]

\end{document}
